\begin{theorem}[全局时钟的消失]\label{thm:conclusion-global-clock-vanishing}
设 \(\mathfrak U=\{U_i\}_{i\in I}\) 为一组区域支撑覆盖，并在每个 \(U_i\) 上给定由局部状态精化诱导的时间读出
\[
\tau_i.
\]
记相应的二阶粘合障碍类为
\[
[\omega(\mathfrak U,\tau)]\in \check H^2(\mathfrak U;A).
\]
若
\[
[\omega(\mathfrak U,\tau)]\neq 0,
\]
则不存在全局可观测量 \(\tau\) 使得
\[
\tau|_{U_i}=\tau_i
\qquad(\forall i\in I).
\]
特别地，不存在观察者无关的全局时钟，也不存在与全部局部证书兼容的统一全局停机泛函。
\end{theorem}
\leanverified{paper\_conclusion\_global\_clock\_vanishing}
\begin{proof}
把 \(\tau_i\) 视为覆盖 \(\mathfrak U\) 上的局部截面。若存在全局 \(\tau\) 满足 \(\tau|_{U_i}=\tau_i\)，则这些局部截面在全部重叠上来自同一全局对象，故按定义 \ref{def:prefix-site-h2-gluing-obstruction} 与命题 \ref{prop:null-as-h2-obstruction}，其二阶粘合障碍类必须为零。此与
\(
[\omega(\mathfrak U,\tau)]\neq 0
\)
矛盾。再由定理 \ref{thm:prefix-site-cech-null-gerbe}，非平凡 Čech-\(H^2\) 障碍正是“局部存在而全局不可中性化”的证书，因此全局时钟与全局终局读出同时被排除。证毕。
\end{proof}

\begin{corollary}[局部有限停机不推出全局有限共识]\label{cor:conclusion-local-finite-halting-no-global-finite-consensus}
\leanverified{paper\_conclusion\_local\_finite\_halting\_no\_global\_finite\_consensus}
在定理 \ref{thm:conclusion-global-clock-vanishing} 的设定下，即使每个局部转导器都具有有限内在停机深度，且每个局部读出都可在有限同步预算内完成，只要
\[
[\omega(\mathfrak U,\tau)]\neq 0,
\]
便不存在一个单一常数 \(T_{\mathrm{glob}}<\infty\)，使全部区域在时间 \(T_{\mathrm{glob}}\) 内达到全局终局共识。
\end{corollary}
\begin{proof}
若这样的 \(T_{\mathrm{glob}}\) 存在，则各 \(U_i\) 上“到达终局”的局部证书可在共同的有限时间截面上同步比较，并组成一组在全部重叠上兼容的局部终局数据。由定理 \ref{thm:conclusion-global-clock-vanishing}，这将迫使相应的全局终局读出存在，与
\(
[\omega(\mathfrak U,\tau)]\neq 0
\)
矛盾。故局部有限停机与局部有限同步并不能推出全局有限共识。证毕。
\end{proof}

\begin{proposition}[区域共识函子的分离而非层化]\label{prop:conclusion-consensus-functor-separated-not-sheaf}
\leanverified{paper\_conclusion\_consensus\_functor\_separated\_not\_sheaf}
对每个区域 \(U\)，记 \(\mathcal C(U)\) 为所有在 \(U\) 上协议可接受且在重叠上可比较的局部共识数据之集合。则函子 \(\mathcal C\) 满足：
\begin{enumerate}
  \item 唯一性成立，因而 \(\mathcal C\) 是分离对象；
  \item 一般不满足有效下降，故 \(\mathcal C\) 未必构成层。
\end{enumerate}
\end{proposition}
\begin{proof}
由定理 \ref{thm:recursive-addressing-readout-separatedness-null}，一旦局部读出在有限覆盖上逐块一致，则相应全局对象若存在必唯一，故 \(\mathcal C\) 具有分离性。另一方面，命题 \ref{prop:null-as-h2-obstruction} 与定理 \ref{thm:prefix-site-cech-null-gerbe} 给出标准情形：局部截面在各 patch 与双重交上都存在并相容，但因非平凡二阶障碍类而不能全球化。故有效下降一般失败，\(\mathcal C\) 不是层。证毕。
\end{proof}

\begin{theorem}[同步 holonomy 的 Beck--Chevalley 证书]\label{thm:conclusion-sync-holonomy-bc-certificate}
设 \(\Sigma\) 为 surjection nerve 中一个有限二维比较曲面，其每个 \(2\)-胞腔 \(\square\subset\Sigma\) 都对应于一个基本 Beck--Chevalley 比较方块
\[
X\xrightarrow{f}Y\xrightarrow{g}Z
\]
以及由同一边界比较数据诱导的局部分布 \(\nu_\square\in\Delta(Z)\)。定义局部同步缺陷
\[
\kappa_\square(\nu_\square)
:=
D_{\mathrm{KL}}\!\bigl(L_fL_g\nu_\square\|L_{g\circ f}\nu_\square\bigr)
=
\mathbb E_{Z\sim \nu_\square}\!\left[\log\frac{\mathrm{AM}_Z}{\mathrm{GM}_Z}\right]\ge 0,
\]
并令
\[
\operatorname{Hol}_{\mathrm{sync}}(\Sigma)
:=
\sum_{\square\subset\Sigma}\kappa_\square(\nu_\square).
\]
则：
\begin{enumerate}
  \item \(\operatorname{Hol}_{\mathrm{sync}}(\Sigma)\ge 0\)；
  \item 在保持同一比较 \(2\)-链与同一边界数据的 chart 细分及 gauge 改写下，\(\operatorname{Hol}_{\mathrm{sync}}(\Sigma)\) 不变；
  \item
  \[
  \operatorname{Hol}_{\mathrm{sync}}(\Sigma)=0
  \iff
  \kappa_\square(\nu_\square)=0\quad(\forall \square\subset\Sigma),
  \]
  等价地，\(\Sigma\) 上每个基本交换方块都严格满足 Beck--Chevalley。
\end{enumerate}
\end{theorem}
\leanverified{paper\_conclusion\_sync\_holonomy\_bc\_certificate}
\begin{proof}
每个 \(\kappa_\square\) 都是 KL 散度，故逐项非负，得 \(\operatorname{Hol}_{\mathrm{sync}}(\Sigma)\ge0\)。对单个方块，定理 \ref{thm:conclusion-discrete-jacobian-strictification-mediated-defect} 已把该缺陷识别为唯一 strictification 核与朴素提升之间的 KL 距离；对由若干方块组成的细分链，定理 \ref{thm:pom-bc-chain-defect-equals-escort-curvature-sum} 把总缺陷精确裂解为逐层 escort 曲率之和，因此在不改变比较 \(2\)-链与边界数据的细分下，总和保持不变。再由命题 \ref{prop:pom-bc-bianchi}，Beck--Chevalley 曲率是一个 \(2\)-上链闭项，允许的 gauge 改写只改变势的代表元而不改变相应的曲率类，因此同一比较曲面的总缺陷不受代表选择影响。

最后，由每一项 \(\kappa_\square\ge0\) 可知总和为零当且仅当逐项为零；而单个 KL 缺陷为零又当且仅当该方块严格 Beck--Chevalley。证毕。
\end{proof}

\begin{corollary}[时间膨胀仅属于比较通道]\label{cor:conclusion-time-dilation-comparison-channel-only}
设 \(T_i\) 为观察者纤维 \(i\) 的内在停机深度。则 \(T_i\) 由局部演化本身确定，不因外部比较而被回写；相反，任一跨 chart 的同步比较所出现的附加“时间膨胀”完全由 \(\operatorname{Hol}_{\mathrm{sync}}\) 度量。特别地，若
\[
\operatorname{Hol}_{\mathrm{sync}}(\Sigma)=0,
\]
则比较通道上不存在额外同步债务；若其为正，则额外延迟纯属比较通道的几何残差，而非内在停机步数的漂移。
\end{corollary}
\begin{proof}
内在停机深度 \(T_i\) 由各观察者纤维上的局部转导与停止规则给出，是局部演化的不变量。任何跨 chart 的额外代价都只能来自比较方块不严格交换时的同步债务，而这正由定理 \ref{thm:conclusion-sync-holonomy-bc-certificate} 中的 \(\operatorname{Hol}_{\mathrm{sync}}\) 量化。故所谓“时间膨胀”不属于局部时间参数本身，而属于比较通道上的附加 holonomy。证毕。
\end{proof}

\begin{theorem}[全局 \texorpdfstring{\(\mathsf{NULL}\)}{NULL} 的局部穷尽与纯上同调剩余]\label{thm:conclusion-global-null-pure-cohomological-remainder}
\leanverified{paper\_conclusion\_global\_null\_pure\_cohomological\_remainder}
设 \(\mathfrak U=\{U_i\}_{i\in I}\) 为一组覆盖，并假定对每个 \(U_i\) 及每个双重交 \(U_{ij}\) 都同时满足：
\begin{enumerate}
  \item 无语义空值；
  \item 无协议空值；
  \item 无碰撞空值。
\end{enumerate}
则下列命题等价：
\begin{enumerate}
  \item \(\bigcup_i U_i\) 上不存在全局可读对象；
  \item 相应的二阶粘合障碍类非平凡。
\end{enumerate}
亦即，一旦三类局部 \(\mathsf{NULL}\) 在 patch 与双重交层面都被穷尽，剩余的全局失败纯粹是 Čech-\(H^2\) 型的上同调残差。
\end{theorem}
\begin{proof}
由定理 \ref{thm:recursive-addressing-readout-separatedness-null}，语义空值、协议空值与碰撞空值分别对应局部存在、下降存在与分离性的失败。题设排除了它们在各 patch 与双重交上的全部出现，因此局部对象存在、一阶兼容性成立且局部读出单值。此时唯一可能残留的全局失败便是二阶粘合障碍。由命题 \ref{prop:null-as-h2-obstruction} 与定理 \ref{thm:prefix-site-cech-null-gerbe}，这种失败当且仅当相应 Čech-\(H^2\) 障碍类非平凡。证毕。
\end{proof}

\begin{proposition}[全局终局态强制 primitive 谱单值化]\label{prop:conclusion-global-terminal-state-forces-global-primitive-spectrum}
设在每个 \(U_i\) 上，局部时间迹通过命题 \ref{prop:time-space-commuting-diagram} 的严格交换接口可重建 primitive 原子谱。若存在一个与全部局部证书兼容的全局终局态，则这些局部 primitive 谱数据必可下降为一个全局单值对象。因而，任一非平凡 Čech 障碍或非零同步 holonomy 都不仅排斥全局时钟，也排斥全局 primitive 谱对象。
\end{proposition}
\begin{proof}
命题 \ref{prop:time-space-commuting-diagram} 把时间侧的迹数据与空间侧的 primitive 原子谱通过同一 \(\zeta\)-接口严格对齐。若全局终局态存在，则局部终局读出形成一组可在覆盖上粘接的全局时间数据；经由该严格交换接口，相应 primitive 谱也必须随之粘接为全局单值对象。于是，一旦定理 \ref{thm:conclusion-global-clock-vanishing} 的二阶障碍非平凡，或定理 \ref{thm:conclusion-sync-holonomy-bc-certificate} 的同步 holonomy 非零而阻断严格比较，就不可能同时得到全局 primitive 谱对象。证毕。
\end{proof}

\begin{theorem}[单一总序时间函子的不可忠实化]\label{thm:conclusion-no-faithful-total-order-time-functor}
记 \(\mathcal O\) 为由观察者纤维局部状态、因果相容更新、重叠比较与共享 forcing 所组成的范畴。若某个有限覆盖上存在非平凡二阶障碍类，则不存在任何忠实函子
\[
F:\mathcal O\longrightarrow (T,\le)
\]
到一个全序集，能够同时保持：
\begin{enumerate}
  \item forcing 的单调性；
  \item 区域比较的可见顺序；
  \item 终局对象的存在性判别。
\end{enumerate}
\end{theorem}
\leanverified{paper\_conclusion\_no\_faithful\_total\_order\_time\_functor}
\begin{proof}
若存在这样的忠实函子 \(F\)，则每个局部状态与其更新链都被压缩为同一全序集上的时间标签，并且区域间比较与终局判别都经由该全序坐标保持。于是局部时间投影在覆盖上被迫拥有一个共同的全局标量代表，从而给出一个观察者无关的全局时钟式对象。此与定理 \ref{thm:conclusion-global-clock-vanishing} 矛盾。故此类 faithful total-order functor 不存在。证毕。
\end{proof}

\begin{conjecture}[零曲率全球化]\label{conj:conclusion-zero-curvature-globalization}
在同一协议类中，若满足：
\begin{enumerate}
  \item 全部局部同步预算有限；
  \item 时间--空间交换接口严格有效；
  \item 所有有限子覆盖上的 Čech-\(H^2\) 障碍类皆平凡；
  \item 全部基本比较方块的 Beck--Chevalley 缺陷恒为零。
\end{enumerate}
则存在一个全局 quasi-time 可观测量与一个全局共识对象。换言之，全局时钟的出现应由二阶平坦性刻画，而不应依赖于任何“实际无穷”极限。
\end{conjecture}
